\subsection{Canonical lattice sizes and lowering by near elements}

Unipotence does not by itself bound how many rational translates can be near a given point. The subspaces of unusually small covolume identify a flag that constrains all such translates.

For a Hermitian lattice \(\Lambda\) at an adelic point,
give each rational subspace and quotient its induced
Euclidean metric at the real places. Write \(f=[E^+:\mathbb Q]\), with \(E^+\) the maximal totally real subfield of \(E\). The size of a nonzero lattice of rank \(j\) is its normalized covolume to the power \(1/(2fj)\). Normalize the real covolume by that of
\(\mathcal O_E^j\) in the standard product metric.

\begin{auxlemma}[The filtration by normalized covolume]
\label{lem:metric-size-filtration}
For an \(\mathcal O_E\)-lattice with the induced subspace and quotient
metrics, there is a unique filtration by saturated sublattices
whose grades are semistable and have strictly increasing sizes.
Here a lattice of size \(t\) is semistable if every nonzero
sublattice has size at least \(t\). If \(F(t)\) contains the grades
of size at most \(t\), then
\[
 \operatorname{size}(Q)\le t < \operatorname{size}(I)
\]
for every nonzero torsion-free quotient \(Q\) of \(F(t)\) and
every nonzero sublattice \(I\) of \(\Lambda/F(t)\), with the two
assertions understood separately. Duality reverses this filtration
and inverts the grade sizes.
\end{auxlemma}
\begin{proof}
\proofstep{Covolumes and attainment.}

For a saturated sublattice \(L'\subset L\) and the quotient
\(L''=L/L'\) with its induced metrics, normalized sizes satisfy
\[
 \operatorname{size}(L)^{\operatorname{rk}L}
 =\operatorname{size}(L')^{\operatorname{rk}L'}
  \operatorname{size}(L'')^{\operatorname{rk}L''}.
\]
This follows by using the determinant norm of the subspace
and the orthogonal projection to the quotient at the
infinite places. The same identity for the integral
determinants holds at every finite place. Saturating a
submodule can only decrease its size.

Among saturated sublattices of fixed rank there are only
finitely many with bounded size. In a fixed determinant ideal
class their exterior line lattices have the form
\(\mathfrak a w\) in the exterior power of \(\Lambda\), for
a fixed ideal representative \(\mathfrak a\). A bound for size
gives a bound for the product of the archimedean norms of \(w\).
Put \(b_v=\log\|w_v\|\) and \(\bar b=f^{-1}\sum_v b_v\).

Choose a unit \(u\) with
\[
     \max_v |\log |u|_v+b_v-\bar b|\le C_E .
\]
Here each complex embedding is taken once. The vector
\((\bar b-b_v)_v\) has sum zero, and \(C_E\) can be chosen as the
radius of a fundamental domain for the logarithms of the units. Dirichlet's unit lattice is
cocompact in the hyperplane of logarithms of product one.
Then \(uw\) generates the same exterior line lattice, since
\(u\) is a unit, and lies in a bounded set of the fixed fractional
exterior lattice and there are only finitely many possibilities.
If \(w\) represents that exterior line, the original
saturated sublattice is recovered by
\[
 L=\{v\in\Lambda:v\wedge w=0\}.
\]
Indeed over \(E\) this kernel is the subspace spanned by the
decomposable exterior vector, and saturation recovers its intersection
with \(\Lambda\).

\proofstep{The canonical increasing filtration.}

Let \(C(L)\) be the normalized covolume of a lattice \(L\),
with \(C(0)=1\) and \(C(L)=\operatorname{size}(L)^{2f\operatorname{rk}L}\).
For two saturated sublattices \(U,V\), put \(I=U\cap V\) and
\(W=(U+V)^{\rm sat}\). Project \(V\) orthogonally off \(U\).
The quotient \(V/I\) maps to the image lattice by a
volume-decreasing map, and saturation can only decrease the
volume further. The determinant formula in the two exact sequences
with quotients \(V/I\) and \(W/U\) therefore gives
\[
 C(I)C(W)\le C(U)C(V),\qquad
 \operatorname{rk}I+\operatorname{rk}W
                  =\operatorname{rk}U+\operatorname{rk}V .
\]
If \(U,V\) have the least size \(t\), put \(d_L=\operatorname{rk}L\).
Then
\[
 t^{2f(d_I+d_W)}\le C(I)C(W)\le C(U)C(V)
              =t^{2f(d_U+d_V)} .
\]
The two outer bounds coincide even when \(I=0\), with \(C(0)=1\).

In particular
\(W\) also has size \(t\). The maximal-rank minimizer is therefore
unique. In the quotient by that minimizer, a sublattice of size at
most \(t\) would have preimage of size at most \(t\) and larger
rank, since the preimage's covolume is the product of the two
factor covolumes. Thus the next step has strictly larger size. This constructs the increasing
Harder--Narasimhan sizes. Each grade of size \(t\) is
semistable, meaning that all its nonzero sublattices have
size at least \(t\). For a grade \(G\) of size \(t\) and a nonzero torsion-free
quotient \(Q=G/K\), multiplicativity and minimality give
\[
 C(Q)=\frac{C(G)}{C(K)}\le
       \frac{t^{2f\operatorname{rk}G}}{t^{2f\operatorname{rk}K}}
       =t^{2f\operatorname{rk}Q}.
\]
This includes \(K=0\).

To prove uniqueness, intersect a nonzero sublattice \(L\) with a
candidate filtration whose grade sizes are \(t_1<\cdots<t_a\).
Let \(d_i\) be the ranks of the induced grades of \(L\).
Each embeds integrally in its ambient grade. Orthogonal projection
to the ambient quotient can only decrease volume, as can saturation.
Multiplicativity on \(L\) therefore gives
\[
 C(L)\ge\prod_{i=1}^a t_i^{2fd_i}\ge t_1^{2f\operatorname{rk}L}.
\]

At equality all \(d_i\) for \(i>1\) vanish. The first step is thus
the maximal-rank minimizer, and induction gives uniqueness on its
quotient. Dualizing reverses the order of the
grades and takes reciprocal sizes, by the determinant
formula for a lattice and its integral dual. The quotient
size property of a semistable grade makes the dual grade
semistable. In particular if \(F(t)\) is the filtration
step with sizes at most \(t\), every nonzero torsion-free
quotient of \(F(t)\) has size at most \(t\), by duality and the
least-size inequality. Every nonzero sublattice of
\(\Lambda/F(t)\) has size strictly greater than \(t\).
\end{proof}

\begin{lemma}[The short flag of a self-dual Hermitian lattice]
\label{lem:canonical-short-flag}
The lattice has a canonical filtration with strictly increasing
sizes and semistable grades. Write \(F(t)\) for its step of sizes
at most \(t\). If the lattice is self-dual then its step \(S\) of
sizes strictly below \(1/2\) satisfies \(S^\perp=F(2)\),
\(\operatorname{rank}S\le q\), and the grades of \(S^\perp/S\)
have sizes in \([1/2,2]\).
If \(a\) is an integral endomorphism, functoriality in the size filtration gives 
\[
 \max_{v\mid\infty}\|a_v\|_{\rm op}\le\delta\quad\Longrightarrow\quad aF(t)\subseteq F(\delta t).
\]
 In particular, if
\(\|\gamma-I\|_{\rm op}\le\delta<1/4\) at all real places,
an integral unitary \(\gamma\) lies in the unipotent radical of
the stabilizer of the short flag. For an empty short flag
this forces \(\gamma=I\).
\end{lemma}

\begin{proof}
Use Lemma~\ref{lem:metric-size-filtration} with the normalized
covolume of the Hermitian lattice.
\proofstep{The perpendicular step and endpoint sizes.}

The Hermitian pairing makes the metric lattice self-dual.
Let \(S\) contain exactly the grades of size strictly less than \(1/2\).
Their dual grades have size strictly greater than \(2\). It follows that
\[
 S\subseteq S^\perp=F(2),\qquad \operatorname{rk}S\le q,\qquad
 \tfrac12\le\operatorname{size}(G)\le2
 \quad(G\text{ a grade of }S^\perp/S).
\] This describes the short flag and
fixes the treatment of the two endpoint sizes.

\begin{construction}
\begin{tikzpicture}[x=1cm,y=1cm,>=Stealth]
 \draw[->] (0,0)--(10.7,0) node[below left] {grade size};
 \draw[line width=.8pt] (3,-.08)--(3,.10) (7,-.08)--(7,.10);
 \node[below] at (3,-.06) {$1/2$};
 \node[below] at (7,-.06) {$2$};
 \node[above] at (1.4,.11) {$S$};
 \node[above] at (5,.11) {$S^\perp/S$};
 \node[above] at (8.9,.11) {$\Lambda/S^\perp$};
 \draw[<->] (1.45,.58) to[bend left=13]
   node[midway,above] {duality \(t\mapsto t^{-1}\)} (8.9,.58);
\end{tikzpicture}
\caption{The short filtration of Lemma~\ref{lem:canonical-short-flag}. Both boundary sizes belong to the middle quotient. The two outer groups of grades are dual.}\label{fig:short-reciprocal-grades}
\end{construction}

\proofstep{Lowering by a near element.}

The canonical filtration is intrinsic to the lattice. It constrains the difference of a near element from the identity, forcing it to lower the appropriate steps.


Suppose an integral endomorphism \(A\) of \(\Lambda\) has
operator norm at most \(\delta\) at every real place.
Let \(A_t:F(t)\to\Lambda/F(\delta t)\) be the induced map.
If its image has rank \(j>0\), pass to the saturated source
quotient \(Q_t\) and put \(I_t=(\operatorname{im}A_t)^{\rm sat}\)
in the target. The complex determinant norm is at most
\(\delta^j\) at each place. With the real squares and the
normalization of \(C\) this gives
\[
 \frac{C(I_t)}{C(Q_t)}\le\delta^{2fj},\qquad
 \operatorname{size}(I_t)\le
            \delta\operatorname{size}(Q_t)\le\delta t .
\]
Saturation accounts for the inequality on the left.
Every nonzero sublattice of the target has size strictly greater
than \(\delta t\), so \(A_t=0\).

Apply this with \(A=\gamma-I\) and the given \(0<\delta<1/4\).
The case \(\delta=0\) gives \(\gamma=I\) directly. It gives identity action
on each short grade, and the image of \(S^\perp=F(2)\)
under \(A\) lies in \(S\). Unitarity gives identity action
on the dual grades as well. Hence \(\gamma\) belongs to the
unipotent radical of the stabilizer of the short flag. If
the short flag is empty, this argument gives \(\gamma=I\).

\end{proof}

\subsection{Packing near elements on the flag blocks}

Lowering the short flag places the near elements in its unipotent radical. To count them one still needs the lattice spacing in each independent block of that radical.

\begin{lemma}[Near elements counted in the short-flag blocks]
\label{lem:near-flag-packing}
Use the grades of Lemma~\ref{lem:canonical-short-flag}.
For integral elements in the unipotent radical of the short flag
whose blocks of \(\gamma-I\) have Hilbert--Schmidt norms at
most \(\delta\) at every real place, the number is bounded by 
\[
 \prod_{v\mid\infty}\prod_{\text{free blocks at }v}(1+12\delta/\lambda)^d.
\]
 Here \(d\) is the block
dimension and
\(\lambda=\tau_i/\tau_j,\tau_i\tau_j,\tau_i,\tau_i^2\)
for the corresponding four block types in \(\mathcal L(\tau)\).
\end{lemma}

\begin{proof}
\proofstep{Induction by block distance.}

Complete the short flag by its perpendicular flag, using
the single middle block \(S^\perp/S\). Choose orthogonal
Euclidean complements at each real place, with dual blocks
paired by the unitary signature matrix. Write the resulting
blocks of \(\gamma-I\) in that order. They are strictly
upper triangular and each has Hilbert--Schmidt norm at
most \(\delta\). Count the possibilities in increasing
distance from the block diagonal. For elements agreeing on
all shorter distances, their difference maps each integral
filtration step into the step lower by the current distance.
Its blocks at that distance are therefore integral maps on
the graded lattices.

For two such elements write \(\gamma_i=1+A_i\) and
\(D=A_1-A_2\). Subtracting the two unitary identities gives
\[
      D+D^\dagger=-(D^\dagger A_1+A_2^\dagger D).
\]
The right side lowers by at least one more step, since
both \(A_i\) are strictly triangular. Thus at the current
distance the adjoint blocks in the difference are negatives
of one another's adjoints. It
suffices to count one block in each such pair and the
self-paired skew-Hermitian differences.

\proofstep{Block separation and volume packing.}

Once the successive free blocks are known, the remaining choices in each block form a separated set. This changes the problem from counting group elements to packing points in Euclidean balls.


For a nonzero difference in one block the same
kernel-to-saturated-image determinant argument bounds the
geometric mean of the operator norms below by the ratio
of the least target size to the largest source size.
For a nonzero block difference \(D\), write
\(l(D)=(\prod_v\|D_v\|_{\rm op})^{1/f}\). The resulting lower bounds are
\[
 l(D)\ge
 \begin{cases}
 \tau_i/\tau_j\ \text{or }\tau_i\tau_j,
                        &\text{the two cross blocks at }i<j,\\
 \tau_i/2,              &\text{a middle block and grade }i,\\
 \tau_i^2,              &\text{the self-paired block at }i .
 \end{cases}
\]
All the relevant Hilbert--Schmidt norms dominate the
operator norms. For an available ratio \(l\), put disjoint
product balls of radius \(l/2\) about the possible blocks
at the real places. An overlap would bound all difference
norms strictly below \(l\). In a self-paired block the
possibilities lie in an affine space with skew-Hermitian
direction, which is the space used for packing. Centering
at any one candidate bounds all radii from that center
by \(2\delta\).

If a finite subset has \(M\) candidates, translation and scaling
of product Lebesgue measure give
\[
 M(l/2)^d\operatorname{vol}(B_1)
       \le(2\delta+l/2)^d\operatorname{vol}(B_1),\qquad
 M\le(1+4\delta/l)^d ,
\]
where \(d\) includes all \(f\) places and \(B_1\) is the unit
product ball in the free block coordinates. This reasoning also
bounds an initially infinite candidate set, by applying the
volume inequality to any finite subset. Iterating through
the block distances determines the full matrices. The
dimensions of the free blocks are those in
\(\mathcal L(\tau)\) in \eqref{eq:count-flag-density-product}. Hence the number is at most
the product over those factors with \(\lambda^d\)
replaced by \((1+12\delta/\lambda)^d\), for
\(\lambda=\tau_i/\tau_j,\tau_i\tau_j,\tau_i,\tau_i^2\).
\end{proof}
